% Bagian: first-order-logic
% Bab: models-theories
% Subbagian: theories

\documentclass[../../../include/open-logic-section]{subfiles}

\begin{document}

\olfileid[id]{fol}{mat}{the}

\olsection{Contoh-Contoh Teori Orde Pertama}

\begin{ex}
Teori urutan linear ketat dalam bahasa~$\Lang L_<$ diaksiomatisasi oleh
himpunan
\begin{align*}
\{\quad & \lforall[x][\lnot x < x], \\
& \lforall[x][\lforall[y][((x < y \lor y <
    x) \lor x = y)]], \\
& \lforall[x][\lforall[y][\lforall[z][((x < y
      \land y < z) \lif x < z)]]] \quad \}
\end{align*}
Himpunan ini sepenuhnya menangkap !!{structure}s yang dimaksudkan: setiap
urutan linear ketat merupakan model sistem aksioma ini, dan sebaliknya, jika
$R$ merupakan urutan linear ketat pada suatu himpunan $X$, maka struktur
$\Struct M$ dengan $\Domain M = X$ dan $\Assign{<}{M} = R$ merupakan model
teori ini.
\end{ex}

\begin{ex}
Teori grup dalam bahasa dengan $\Obj 1$ (!!{constant}) dan $\cdot$
(!!{function} dua-tempat) diaksiomatisasi oleh
\begin{align*}
& \lforall[x][\eq[(x \cdot \Obj 1)][x]]\\
& \lforall[x][\lforall[y][\lforall[z][\eq[(x \cdot (y \cdot z))][((x
          \cdot y) \cdot z)]]]]\\
& \lforall[x][\lexists[y][\eq[(x \cdot y)][\Obj 1]]]
\end{align*}
\end{ex}

\begin{ex}
Teori aritmetika Peano diaksiomatisasi oleh kalimat-kalimat berikut dalam
bahasa aritmetika~$\Lang L_A$.
\begin{align*}
& \lforall[x][\lforall[y][(\eq[x'][y'] \lif \eq[x][y])]]\\
& \lforall[x][\eq/[\Obj 0][x']]\\
& \lforall[x][\eq[(x + \Obj 0)][x]]\\
& \lforall[x][\lforall[y][\eq[(x + y')][(x + y)']]]\\
& \lforall[x][\eq[(x \times \Obj 0)][\Obj 0]]\\
& \lforall[x][\lforall[y][\eq[(x \times y')][((x \times y) + x)]]]\\
& \lforall[x][\lforall[y][(x < y \liff \lexists[z][\eq[(z' + x)][y])]]]\\
\intertext{ditambah semua kalimat berbentuk}
& (!A(\Obj 0) \land \lforall[x][(!A(x) \lif !A(x'))]) \lif \lforall[x][!A(x)]
\end{align*}
Karena terdapat tak terhingga banyak kalimat berbentuk seperti yang terakhir,
sistem aksioma ini tak hingga. Bentuk terakhir tersebut disebut
\emph{skema induksi}. (Sebenarnya, skema induksi sedikit lebih rumit daripada
yang kita paparkan di sini.)

Aksioma terakhir merupakan \emph{definisi eksplisit} dari~$<$.
\end{ex}

\begin{ex}
Teori himpunan murni memegang peranan penting dalam dasar-dasar (dan dalam
filsafat) matematika. Suatu himpunan disebut murni jika semua !!{element}snya juga
merupakan himpunan murni. Oleh sebab itu, himpunan kosong tergolong murni,
tetapi suatu himpunan yang mempunyai sesuatu yang bukan himpunan sebagai
!!a{element} tidaklah murni. Jadi, himpunan murni ialah himpunan yang dibentuk
hanya dari himpunan kosong dan tanpa ``urelemen,'' yaitu objek yang bukan
merupakan himpunan.

Berikut ini dapat dipandang sebagai suatu sistem aksioma bagi teori himpunan
murni:
\begin{align*}
& \lexists[x][\lnot \lexists[y][y \in x]]\\
& \lforall[x][\lforall[y][(\lforall[z](z \in x \liff z \in y) \lif
\eq[x][y])]]\\
& \lforall[x][\lforall[y][\lexists[z][\lforall[u][(u \in z \liff
(\eq[u][x] \lor \eq[u][y]))]]]]\\
& \lforall[x][\lexists[y][\lforall[z][(z \in y \liff \lexists[u][(z \in
        u \land u \in x)])]]]\\
\intertext{ditambah semua kalimat berbentuk} &
\lexists[x][\lforall[y][(y \in x \liff !A(y))]]
\end{align*}
Aksioma pertama menyatakan bahwa terdapat suatu himpunan tanpa !!{element}s
(yakni, $\emptyset$ ada); aksioma kedua menyatakan bahwa himpunan bersifat
ekstensional; aksioma ketiga bahwa untuk setiap himpunan $X$ dan $Y$, himpunan
$\{X, Y\}$ ada; dan aksioma keempat bahwa untuk setiap himpunan $X$, himpunan
$\cup X$ ada, dengan $\cup X$ sebagai gabungan semua anggota~$X$.

!!^{sentence}s yang disebutkan terakhir secara bersama-sama dinamakan
\emph{skema komprehensi naif}. Pada dasarnya, skema itu menyatakan bahwa bagi
setiap $!A(x)$, himpunan $\Setabs{x}{!A(x)}$ ada---yang sepintas tampak sebagai
aksioma yang benar, berguna, dan bahkan mungkin diperlukan. Skema itu disebut
``naif'' karena ternyata membuat teori ini tak dapat dipenuhi: jika kita
mengambil $!A(y)$ sebagai $\lnot y \in y$, kita memperoleh !!{sentence}
\[
\lexists[x][\lforall[y][(y \in x \liff \lnot y \in y)]]
\]
dan !!{sentence} ini tidak terpenuhi dalam !!{structure} mana pun.
\end{ex}

\begin{ex}
Dalam bidang \emph{mereologi}, relasi \emph{bagian-dari} merupakan relasi
mendasar. Sama seperti teori himpunan, terdapat teori-teori bagian-dari yang
mengaksiomatisasi beragam konsepsi (yang kadang saling bertentangan) tentang
relasi ini.

Bahasa mereologi memuat satu simbol predikat dua-tempat~$\Obj P$, dan
$\Atom{\Obj P}{x, y}$ ``berarti'' bahwa $x$ merupakan bagian dari~$y$. Apabila
kita mengacu pada interpretasi ini, !!a{structure} bagi bahasa tersebut
disebut \emph{struktur bagian-dari}. Tentu saja, tidak setiap struktur bagi
satu predikat dua-tempat benar-benar layak disebut demikian. Agar mungkin
menangkap ``bagian-dari,'' $\Assign{\Obj P}{M}$ harus memenuhi beberapa syarat,
yang dapat kita tetapkan sebagai aksioma bagi teori bagian-dari. Misalnya,
bagian-dari merupakan urutan parsial pada objek: setiap objek merupakan bagian
(meskipun bagian \emph{taksejati}) dari dirinya sendiri; tidak ada dua objek
berbeda yang masing-masing merupakan bagian dari yang lain; dan bagian dari
suatu bagian dari sebuah objek juga merupakan bagian dari objek tersebut.
Perhatikan bahwa dalam pengertian ini ``merupakan bagian dari'' menyerupai
``merupakan himpunan bagian dari,'' tetapi tidak menyerupai ``merupakan anggota
dari,'' yang tidak refleksif maupun transitif.
\begin{align*}
& \lforall[x][\Atom{\Obj P}{x,x}] \\
& \lforall[x][\lforall[y][((\Part{x}{y} \land \Part{y}{x})
      \lif \eq[x][y])]] \\
& \lforall[x][\lforall[y][\lforall[z][((\Part{x}{y} \land
        \Part{y}{z}) \lif \Part{x}{z})]]]\\
\intertext{Selain itu, setiap dua objek mempunyai suatu jumlah mereologis
  (sebuah objek yang mempunyai kedua objek itu sebagai bagian dan minimal
  dalam hal ini).}  &
\lforall[x][\lforall[y][\lexists[z][\lforall[u][(\Part{z}{u} \liff
        (\Part{x}{u} \land \Part{y}{u}))]]]]
\end{align*}
Ini hanyalah sebagian dari prinsip-prinsip dasar bagian-dari yang dibahas para
metafisikawan. Namun, prinsip-prinsip lanjutan dengan cepat menjadi sukar
dirumuskan atau dituliskan tanpa terlebih dahulu memperkenalkan sejumlah
relasi terdefinisi. Misalnya, kebanyakan metafisikawan yang tertarik pada
mereologi juga memandang prinsip berikut sebagai sahih: setiap kali suatu
objek~$x$ mempunyai bagian sejati~$y$, objek itu juga mempunyai suatu
bagian~$z$ yang tidak mempunyai bagian bersama dengan~$y$, dan sedemikian
sehingga fusi $y$ dan $z$ adalah~$x$.
\end{ex}

\end{document}
